\section{Discussion}\label{sec:disc}

\subsection{What the chain can and cannot tell a lender}
Chain data is strong on liabilities and cash. It gives an exact loan book, every borrower address, every payment, every holder and every queued exit, and it matches the vault's own accounting to within a fraction of a percent. It is weak on assets. It cannot show that a GPU exists, that it is rented, at what price, or that a borrower earns enough to pay. Three of our findings sit at that boundary. First, more of the cash yield has come from stablecoin issuers (M\textsuperscript{0} and PYUSD) than from borrowers, so holders' returns depend on those programs as much as on GPU economics. Second, \DTShare{} of assets is an escrow claim on funds that left the protocol's contracts before any collateral was posted. Third, the loans run \YearsToMaturity{} more years on average while holders can queue to exit in two weeks, so the vault relies on new deposits and idle cash to pay exits. The contracts permit all three, and none shows up in the headline NAV.

\subsection{What a neocloud unwind would do}\label{sec:unwind}
The borrowers are GPU operators, the class of company now called neoclouds. That sector carries more than \$20\,B of GPU-backed debt from lenders such as Blackstone, PIMCO, Carlyle and Macquarie~\cite{friedman2025}, and its stress is visible: CoreWeave's five-year credit default swaps traded at distressed levels in October 2026~\cite{gokhshtein2026}. H100 rental prices fell from \$7 to \$10 per hour in early 2024 to \$2 to \$4 by late 2025~\cite{friedman2025}, and the GetDeploying series in this paper puts the centralized median at \CentralH{}. The debt structures assume GPUs keep enough value over five to six years to support the loans, while some operators and investors argue the economic life is two to three years~\cite{friedman2025}.

USD.AI is a small on-chain slice of the same trade, with the same collateral and the same borrowers' customers. Table~\ref{tab:stress} shows what defaults would do to the vault, using the reconciled state at block \StateBlock{} and a recovery rate on the GPU collateral. This is arithmetic, not a credit model; it ignores the time to repossess and sell hardware, legal costs and the fact that recoveries fall when every lender sells GPUs at once. The largest borrower holds \TopOneOut{} of outstanding principal and the three largest \TopThreeOut{}. At a 40\% recovery, a default by the largest borrower costs \StressOnePct{} of assets and leaves NAV at \StressOneNAV{}; the three largest cost \StressThreePct{} and push NAV to \StressThreeNAV{}, below par. If every borrower defaults at the same recovery, NAV falls to \StressAllNAV{}. If the open escrow is also not returned, the loss is \StressWorstPct{} of assets and NAV is \StressWorstNAV{}.

\begin{table}[t]
\caption{Loss to vault assets under default scenarios (arithmetic on the reconciled state; recovery is on GPU collateral)}
\label{tab:stress}
\centering
\scriptsize
\resizebox{\columnwidth}{!}{\begin{tabular}{@{}llrrr@{}}
\toprule
Defaults & Recovery & Loss & Share of assets & NAV after \\
\midrule
largest borrower & 60\% & \$21.8\,M & 3.9\% & 1.076 \\
three largest borrowers & 60\% & \$51.4\,M & 9.3\% & 1.017 \\
every borrower & 60\% & \$72.9\,M & 13.1\% & 0.973 \\
largest borrower & 40\% & \$32.6\,M & 5.9\% & 1.054 \\
three largest borrowers & 40\% & \$77.1\,M & 13.9\% & 0.965 \\
every borrower & 40\% & \$109.3\,M & 19.7\% & 0.900 \\
largest borrower & 20\% & \$43.5\,M & 7.8\% & 1.033 \\
three largest borrowers & 20\% & \$102.8\,M & 18.5\% & 0.913 \\
every borrower & 20\% & \$145.8\,M & 26.3\% & 0.826 \\
every borrower; escrow not returned & 40\% & \$238.2\,M & 42.9\% & 0.640 \\
\bottomrule
\end{tabular}
}
\end{table}

Three features of the design decide how such a loss would play out. First, the vault prices deposits at an optimistic NAV and redemptions at a conservative one, and nothing in the contracts marks a loan down before a breach event. A borrower that stops paying would show up as a missed monthly repayment, then as a \texttt{LoanBreached} event, and only then in NAV. Holders who read the queue could leave at par before the mark. Second, exits are a queue. Today the idle deposit balance covers the pending queue \IdleOverPending{} times over, so the vault can pay out without touching a loan. In a run the idle balance goes first, then the vault has to wait on \YearsToMaturity{} more years of loan maturities. Third, sUSDai is collateral elsewhere: \EthMorpho{} of the Ethereum supply is in Morpho Blue and \PendleArb{} of the Arbitrum supply is in Pendle. A NAV mark below par would trigger liquidations in those venues that sell sUSDai into thin pools, pushing the market price further below NAV and feeding the queue.

The same mechanics apply to every protocol that lends stablecoin deposits against GPUs, including those that issue their own synthetic dollar~\cite{gaib2025}. An unwind would hit them together, because they share customers (AI labs renting compute), collateral (the same GPU generations), and the same secondary market for used hardware. Three questions separate the ones that would survive: how much of the yield comes from borrowers rather than from stablecoin incentives; how many borrowers there are and how large the biggest one is; and whether redemptions are backed by idle cash or by the hope of new deposits. All three can be measured from the chain with the code in this repository, and none of them appears in a headline APY.

\subsection{Risks the chain shows only in outline}
Seven risks are visible on-chain but cannot be sized from it.

\emph{Related parties.} The same accounts submit the protocol's multisig transactions and the escrow facility's. The largest borrower, the fee recipient and the escrow facility all route PYUSD to the same off-ramp wallet. None of this is improper on its face, but a lender pricing ``\NBorrowers{} independent borrowers'' needs the entity map behind those wallets, and the chain does not supply it.

\emph{Key compromise.} The admin role can pause every core contract, and the escrow and multisig are controlled by a handful of externally owned accounts. Two days before this paper was finished, a leak of environment variables at a hosting provider let one attacker drain 231 wallets on 17 chains in about an hour, because the private keys sat on application servers~\cite{railway2026}. The same failure mode applied to a signer here would not drain the vault directly, but it could pause redemptions or move escrow.

\emph{Fixed loans, floating benchmark.} The loans are fixed at \RateMin{} to \RateMax{} for \YearsToMaturity{} more years. The yield holders compare against moves: Treasury bills, and the PYUSD and M\textsuperscript{0} incentives that supplied most of the cash yield so far. If those fall, the vault's spread widens; if they rise, the vault holds below-market paper and the queue fills.

\emph{Custody, title and valuation.} The hardware is described on-chain down to the server model and, for the early loans, the serial number (Section~\ref{sec:collateral}). What is not on-chain is who holds it, where, under what security filing or insurance, and whether the stated collateral values are appraisals or purchase prices. At \BookLTV{} book LTV on a three-year straight-line life, the margin for error in those values is small.

\emph{Bridged exits.} \BridgedShare{} of shares live on Ethereum, Base, Plasma and Solana. A holder there must bridge back to Arbitrum before joining the queue. A paused bridge pauses the exit for two fifths of the supply.

\emph{Leverage on top.} Pendle holds \PendleArb{} of Arbitrum sUSDai and Morpho Blue \EthMorpho{} of Ethereum sUSDai. Fixed-yield buyers in Pendle have locked in today's rate; a NAV mark hits them first, and the leveraged positions sit there.

\emph{Classification.} This is pooled private credit with a yield-bearing token. Payment-stablecoin rules in the United States now distinguish tokens that pay yield from those that do not, and the treatment of a structure like this one is not settled. A ruling could change who may hold it or how it may be marketed.

\subsection{What this means for different readers}
\emph{Traders.} Three on-chain signals lead a NAV mark: a missed monthly repayment (today the longest gap is \MaxDaysSincePayment{} days), a rise in the queue relative to idle cash (today \IdleOverPending{}$\times$ covered), and the DEX discount of sUSDai to NAV. All three are computed by the repository and can be polled. The conservative redemption price (\RedemptionPrice{}) against the deposit price (\DepositPrice{}) is the protocol's own view of how much accrued interest is at risk.

\emph{Institutional allocators.} Underwrite the escrow before the loans: \DTShare{} of assets is a claim with no collateral yet, most likely a refinancing in progress. Price the yield on borrower cash alone (\KappaBorrowers{} coverage), not the headline. Treat the loan book as one neocloud exposure with \TopOneOut{} in a single name at \BookLTV{} LTV on three-year-life hardware, and stress it against the recovery table (Table~\ref{tab:stress}). Ask who appraised the stated collateral values and what the fleet's utilization is; the chain gives the model and count but not the revenue.

\emph{Neoclouds.} On-chain credit is cheaper to audit than bank credit, and USD.AI already publishes model, count and stated value for every pledge. A borrower that added utilization, rental revenue and its power contract would be the first a lender could underwrite from public data alone, and could reasonably ask for a lower rate than the \RateWavg{} weighted average, which prices what is still unknown.

\emph{Researchers.} The dataset is frozen and hash-verified, every number reproduces offline, and the method transfers to any GPU-credit protocol on an indexed chain. The open questions are the ones the chain cannot answer: independent collateral valuations, borrower revenue, and recoveries once the first loan defaults. A second protocol and a longer window would make this a comparative study.

\subsection{Monitoring indicators}
The pipeline can run on a schedule. We suggest six indicators, each computed by code in the repository:
\begin{enumerate}
\item \emph{Cash coverage} $\kappa$ by month, split into borrower interest and base yield.
\item \emph{Pre-funded share}: \texttt{depositTimelockBalance} over \texttt{totalAssets}, with the age of each open escrow.
\item \emph{Borrower HHI} on outstanding principal, and the largest borrower's share.
\item \emph{Queue pressure}: pending redemptions against idle deposit balance, and the age of the oldest pending request.
\item \emph{Missed payments}: loans whose last repayment is older than their repayment interval.
\item \emph{Peg stress}: the market discount of USDai and sUSDai to NAV.
\end{enumerate}

\subsection{Other GPU-backed protocols}
Other protocols also tokenize GPU financing or GPU revenue, and some issue their own synthetic dollars, such as GAIB's AID~\cite{gaib2025}. The same method works for any of them whose contracts run on a chain GoldRush indexes: decode from source, rebuild the loan book, reconcile against state and trace the cash. Comparing them on cash coverage, pre-funding and borrower concentration is a natural next study.

\subsection{Data coverage}
Without a unified API, several steps would have needed an archive node and custom indexing. GoldRush provided complete decoded histories for the lending contracts, event scans across 182 block windows, holder lists on four chains, holder lists at past blocks, one call listing every chain an address has used, cross-chain wallet balances and daily token prices. Three gaps remain. GoldRush does not offer \texttt{eth\_call}, so the state check used a public RPC. Holder lists at past blocks are available only on its Foundational chains, so we have them for Ethereum but not Arbitrum. Akash and other Cosmos chains are not indexed.

\subsection{Limitations}
We study one protocol over 17 months. The loan book is rebuilt from events and depends on our decoder. We check it against vault state, but only at one block. Two topic hashes on the V1 router match no signature in the source history. Current holder lists give concentration at a single point in time. Contract labels come from GoldRush metadata and can be wrong. The GPU price series is a weekly median across clouds, not the price any borrower actually receives. Livepeer payments mix video transcoding and AI inference, and the Akash figures come from the network's own API. We make no claim about any off-chain entity. Wallets are described only by their on-chain role, and nothing here shows that funds were misused.

\subsection{Ethics and right of reply}
All data is public. We do not name individuals or link wallets to companies. We plan to share the escrow findings with USD.AI before publication and to include any response in a revised version.

\section{Conclusion}
Chain data alone was enough to rebuild this GPU-backed dollar in detail. It shows \Originated{} of GPU loans to \NBorrowers{} borrowers, a vault whose cash yield relies on stablecoin issuers, a redemption queue cleared in large monthly batches, and \DTShare{} of assets held as pre-funded escrow instead of collateralized loans. It also shows that on-chain spending on compute is small next to on-chain lending for compute. Until GPU usage and revenue are recorded on-chain, lenders and curators have to estimate that part of the risk from off-chain information.

\section*{Disclosures}
\textbf{Conflicts of interest:} The author has no financial relationship with USD.AI, GAIB, or any protocol, borrower or lender discussed in this paper, and holds no position in their tokens. \textbf{Use of AI:} An AI assistant (Claude) helped write the code and draft the text. The author reviewed the method and results, and the published code generates every number. \textbf{Data:} GoldRush API by Covalent; Akash Console API; GetDeploying GPU price snapshot (CC~BY~4.0)~\cite{getdeploying}; Arbitrum and Solana public RPCs.
